% Generated draft: requires mathematical and source review.
\documentclass[11pt]{article}
\usepackage[margin=25mm]{geometry}
\usepackage{fontspec}
\setmainfont{Latin Modern Roman}
\usepackage{amsmath,amssymb,mathtools,unicode-math}
\setmathfont{Latin Modern Math}
\usepackage{enumitem,needspace,xurl,hyperref}
\setlength{\parindent}{0pt}
\setlength{\parskip}{5pt}
\setlength{\emergencystretch}{3em}
\newcommand{\N}{\mathbb{N}}
\newcommand{\Z}{\mathbb{Z}}
\newcommand{\Q}{\mathbb{Q}}
\newcommand{\R}{\mathbb{R}}
\newcommand{\C}{\mathbb{C}}
\newcommand{\F}{\mathbb{F}}
\newcommand{\A}{\mathbb{A}}
\newcommand{\PP}{\mathbb{P}}
\newcommand{\E}{\mathbb{E}}
\newcommand{\eps}{\varepsilon}
\newcommand{\dd}{\,\mathrm{d}}
\begin{document}
\section*{UnsolvedMath ID: 800259}
\textbf{Source problem code:} catalogue entry
\subsection{Definitions and mathematical statement}
For topological spaces \(X\) and \(Y\), a subset \(U\subseteq X\), and a map \(f:U\to Y\), a continuous extension to the closure means a continuous map \(\widetilde f:\overline U^{\,X}\to Y\) satisfying \(\widetilde f|_U=f\). This is a general definition only; the supplied record does not specify these spaces or the hypotheses on \(f\).

The supplied question is: Does \(f\) have a continuous extension to \(\overline U\)? The source fragment adds the qualification that this is true if \(U\) and \(V\) are pseudoconvex.

The record does not introduce \(U\), \(V\), or \(f\), specify their dimensions or ambient spaces, state any boundary assumptions, or identify the target space of the extension. Consequently, neither the intended quantified problem nor the hypotheses underlying the parenthetical assertion can be recovered precisely. The parenthetical assertion is preserved as a source qualification, not asserted here as a theorem under the general definition above.
\subsection{Short English statement}
The source asks whether an unspecified map \(f\) extends continuously to the closure of \(U\), and states that the answer is affirmative when \(U\) and \(V\) are pseudoconvex. The missing setup prevents a precise formulation of the intended problem.
\subsection{Sources}
\begin{itemize}
\item Supplied local UnsolvedMath registry record: permanent ID \(800259\), exact source code \texttt{catalogue entry}, title ``Several Complex Variables Problem List --- Problem 58.'' The registry attributes the problem to Buzzard, Fornaess, and Krantz, with proposed year 1995, and labels it open; this status is supplied curation, not an independently verified assessment.
\item Buzzard, Fornaess, and Krantz, \emph{The Several Complex Variables Problem List} (1995), Problem 58. Supplied source URL: \url{https://arxiv.org/abs/math/9509202}. The record reports access on 2026-10-16 and extraction by \texttt{arxiv\_latex}; these are supplied provenance details, not a claim of independent consultation.
\end{itemize}
\end{document}
